\documentclass[12pt]{article}
\usepackage[margin=1in]{geometry}
\usepackage{parskip}
\usepackage{xcolor}
\usepackage{soul}
\usepackage{amssymb}

\definecolor{DukeBlue}{RGB}{0,0,156}

\title{Enhanced Presentation Script with Timestamps:\\The Development and Outlook of AI Techniques in Neuroscience}
\author{Juntang Wang}
\date{COMPSCI 402 Final Presentation\\October 13, 2025}

\begin{document}

\maketitle

\section*{Instructions}
This enhanced script covers 15 content slides in **16-18 minutes**. Timestamps show when each slide should start. Read naturally with enthusiasm and use examples to connect with your audience.

\textbf{Key Symbols:}
\begin{itemize}
  \item \textcolor{DukeBlue}{\textbf{[GESTURE]}} / \textcolor{DukeBlue}{\textbf{[POINT]}} = Physical actions
  \item \textcolor{red}{\textbf{[PAUSE]}} = Brief pause for emphasis
  \item \textbf{\ul{Underlined}} = Emphasize this word/phrase
  \item \textcolor{green!60!black}{\textbf{[Example]}} = Real-world context
\end{itemize}

\hrulefill

\section*{Slide 1: Title Slide [0:00-0:15]}

\textcolor{red}{\textbf{[PAUSE 2-3 seconds after slide appears]}}

\textcolor{DukeBlue}{\textbf{[SMILE, look at camera]}}

Hi everyone. My name is Juntang Wang, and I'll present my final report on the development and outlook of AI techniques in neuroscience, completed for the COMPSCI 402 course project.

\hrulefill

\section*{Slide 2: Overview [0:15-1:15]}

\textcolor{DukeBlue}{\textbf{[GESTURE to slide]}}

Today we'll explore a fascinating intersection of two fields. Let me give you a roadmap. First, I'll show you the \ul{bidirectional relationship} between AI and neuroscience—how they help each other. Then, before diving deep, I'll clarify key concepts to make sure we're all on the same page. 

\textcolor{red}{\textbf{[PAUSE]}}

After that, we'll see how this field evolved from simple perceptrons in the 1950s to today's sophisticated systems. The core of my presentation is a \ul{taxonomy}—a systematic organization of AI applications in neuroscience across three categories. We'll examine representative methods, look at impressive performance results, and then discuss the challenges that still keep researchers up at night. Finally, I'll conclude with what this all means for the future.

\hrulefill

\section*{Slide 3: Motivation [1:15-2:30]}

\textcolor{DukeBlue}{\textbf{[LOOK AT CAMERA, enthusiastic]}}

So why should we care about AI in neuroscience? \textcolor{red}{\textbf{[PAUSE]}} The key insight is this is a \ul{two-way street}.

\textcolor{DukeBlue}{\textbf{[GESTURE left]}} On one hand, AI gives neuroscience powerful tools. Think about it—we can now decode what someone is thinking just from brain scans, analyze complex neural signals that humans couldn't interpret, build brain-computer interfaces that let paralyzed patients control robotic arms, and diagnose Alzheimer's years before symptoms appear.

\textcolor{DukeBlue}{\textbf{[GESTURE right]}} On the other hand, neuroscience inspires AI innovation. \textcolor{green!60!black}{\textbf{[Example]}} For instance, the entire field of deep learning was inspired by how our visual cortex processes images in layers. Today, researchers are building neuromorphic chips that mimic brain neurons and use a fraction of the power of traditional computers.

\textcolor{DukeBlue}{\textbf{[POINT to diagram]}} This diagram captures it beautifully—AI provides tools while neuroscience provides inspiration. \textcolor{red}{\textbf{[PAUSE]}} This \ul{synergy} is what makes this field so exciting and productive.

\hrulefill

\section*{Slide 4: Key Concepts [2:30-3:45]}

\textcolor{DukeBlue}{\textbf{[GESTURE to slide]}}

Before we go further, let me make sure we're all on the same page with some key terms. I know audiences come from different backgrounds, so let's quickly clarify these concepts.

\textcolor{DukeBlue}{\textbf{[POINT to left column]}} First, the core terms. When I say \ul{AI or Artificial Intelligence}, I'm talking about systems that learn from data to make decisions—like recognizing patterns in brain scans. \ul{Neuroscience} is the study of the nervous system and brain. And \ul{BCI}—brain-computer interface—is a direct communication pathway between your brain and external devices. \textcolor{green!60!black}{\textbf{[Example]}} Imagine someone who's paralyzed being able to type on a computer just by thinking about moving their hand.

\textcolor{DukeBlue}{\textbf{[POINT to right column]}} Now, the recording methods. \ul{EEG}, or electroencephalography, records electrical brain activity through electrodes on your scalp—it's non-invasive and has great time resolution. \ul{fMRI}, functional MRI, images brain activity by detecting blood flow changes—it gives us beautiful spatial detail. And \ul{neuromorphic computing} refers to hardware chips that actually mimic how biological neurons work, using spikes of electricity like real brains do.

\textcolor{red}{\textbf{[PAUSE]}} These terms will come up throughout, so now we're all speaking the same language.

\hrulefill

\section*{Slide 5: Three Directions \& Challenges [3:45-5:15]}

\textcolor{DukeBlue}{\textbf{[GESTURE to slide]}}

From my literature review, AI in neuroscience falls into \ul{three main directions}.

\textcolor{DukeBlue}{\textbf{[POINT left, count on fingers]}} \ul{First}, neural decoding and brain-computer interfaces—extracting meaningful information from brain signals. \textcolor{green!60!black}{\textbf{[Example]}} This includes work where researchers reconstructed what movie someone was watching just from their brain activity.

\ul{Second}, disease diagnosis through AI-powered neuroimaging. Deep learning models can now detect Alzheimer's with over ninety-eight percent accuracy, often catching it years before traditional methods.

\ul{Third}, neuroscience-inspired AI architectures—systems like spiking neural networks and neuromorphic chips that work more like actual brains.

\textcolor{red}{\textbf{[PAUSE, shift to serious tone]}}

\textcolor{DukeBlue}{\textbf{[POINT right]}} But this field isn't all success stories. We face five major challenges. Data \ul{complexity}—brain data spans from individual neurons to networks of billions. \ul{Multimodal integration}—combining EEG, fMRI, and behavioral data is incredibly difficult. \ul{Interpretability}—doctors won't trust a "black box" AI they can't understand. \ul{Generalizability}—what works for one person often fails for another due to individual brain differences. And \ul{ethics}—if we can read brain signals, how do we protect mental privacy?

\textcolor{red}{\textbf{[PAUSE]}} My report focuses on tackling these challenges, emphasizing work from the \ul{past five years}.

\hrulefill

\section*{Slide 6: Evolution \& Current State [5:15-6:45]}

\textcolor{DukeBlue}{\textbf{[GESTURE to timeline]}}

Let me give you some historical context. \textcolor{DukeBlue}{\textbf{[TRACE timeline left to right]}} This field has come a long way.

It started in the 1950s with the perceptron—a simple attempt to model how neurons work. Not much happened until the 1980s when backpropagation gave neural networks the ability to actually learn. The 1990s and 2000s brought CNNs and RNNs, directly inspired by how our visual cortex is organized in layers.

Then came \ul{2012}—the year everything changed. AlexNet won the ImageNet competition by a huge margin using deep learning, and suddenly everyone realized these brain-inspired networks could do things we thought required human intelligence. Today we have transformers powering ChatGPT and neuromorphic chips that process information like biological brains.

\textcolor{DukeBlue}{\textbf{[GESTURE to columns]}} The achievements are impressive—we can reconstruct visual experiences from brain scans, diagnose diseases at expert level, and build chips that are incredibly energy efficient. \textcolor{red}{\textbf{[PAUSE]}} But gaps remain: no unified brain models, limited interpretability, poor generalization across individuals, and we still need better ethical frameworks.

\hrulefill

\section*{Slide 7: Taxonomy Overview [6:45-7:45]}

\textcolor{DukeBlue}{\textbf{[GESTURE to diagram with enthusiasm]}}

Here's the heart of my report—a taxonomy that organizes this complex field.

\textcolor{DukeBlue}{\textbf{[TRACE tree structure]}} I've organized AI techniques in neuroscience into three primary categories. \ul{First}, neural decoding and BCIs—with approaches using EEG, fMRI, and spiking neural networks. \ul{Second}, disease diagnosis—including Alzheimer's detection, general neuroimaging, and multimodal data integration. \ul{Third}, neuroscience-inspired AI—with deep networks as brain models, neuromorphic computing, and network neuroscience approaches.

\textcolor{red}{\textbf{[PAUSE]}} Think of this as a map. Each branch represents a different way AI and neuroscience intersect. Now let's explore each category in detail.

\hrulefill

\section*{Slide 8: Neural Decoding [7:45-8:45]}

\textcolor{DukeBlue}{\textbf{[GESTURE to slide]}}

Starting with neural decoding—reading information from the brain. For EEG-based approaches, \ul{EEGNet} is a star performer. It's a compact convolutional network that works across multiple BCI paradigms. \textcolor{green!60!black}{\textbf{[Example]}} Whether you're detecting P300 responses when someone sees a target image, or decoding motor intentions for robotic control, EEGNet handles it elegantly with fewer parameters than competing methods.

For neuroimaging, deep networks can now reconstruct visual experiences from fMRI data. Imagine being able to see what someone saw in a movie just by looking at their brain activity—that's what motion-energy encoding models achieve by cleverly separating fast visual information from slow blood flow responses.

And then there's spiking neural networks—SuperSpike and LSTM-SNNs—which are more biologically plausible because they use discrete spikes like real neurons rather than continuous activations. This matters for learning-to-learn scenarios where systems need to adapt quickly, just like biological brains do.

\hrulefill

\section*{Slide 9: Disease Diagnosis [8:45-9:45]}

\textcolor{DukeBlue}{\textbf{[POINT to slide, excited]}}

Now disease diagnosis—this is where AI is already changing clinical practice.

\ul{AlzheimerNet} achieves \textcolor{DukeBlue}{\textbf{ninety-eight point six seven percent accuracy}} classifying all five stages of Alzheimer's plus normal controls. \textcolor{red}{\textbf{[PAUSE]}} That's remarkable. It uses data augmentation and CLAHE image enhancement to handle unbalanced datasets—a common problem in medical imaging where you have way more normal scans than diseased ones.

\textcolor{green!60!black}{\textbf{[Example]}} More recently, researchers adapted YOLOv11—originally designed for finding objects in photos—to detect Alzheimer's by fusing MRI and DTI images. It achieved ninety-three point six percent precision and ninety-six point seven percent mean average precision. This shows how techniques from computer vision can be creatively adapted for medical applications.

Broadly, CNNs have become the methodology of choice across classification, detection, and segmentation tasks. The integration of multimodal data—combining different types of brain scans with behavioral measures—is pushing diagnostic capabilities even further.

\hrulefill

\section*{Slide 10: Neuroscience-Inspired AI [9:45-10:45]}

\textcolor{DukeBlue}{\textbf{[GESTURE to slide]}}

Now we flip the direction—where neuroscience gives back to AI.

Deep neural networks aren't just useful tools; they're also surprisingly good \ul{models of how brains work}. Studies show DNNs can model neural responses in the ventral visual pathway, revealing gradients of feature complexity from simple edges to complex objects. Remarkably, these representations are shared across different people, suggesting common computational principles.

\textcolor{DukeBlue}{\textbf{[POINT with emphasis]}} But the really exciting development is neuromorphic hardware. \ul{Intel's Loihi processor}—a neuromorphic manycore chip—shows \textcolor{DukeBlue}{\textbf{over one thousand times}} better energy efficiency than conventional CPUs for certain tasks. \textcolor{red}{\textbf{[PAUSE]}} Think about what that means. \textcolor{green!60!black}{\textbf{[Example]}} Your brain runs on about twenty watts—less than a lightbulb—yet outperforms supercomputers on many tasks. Neuromorphic chips try to capture that efficiency by using spike-based computation like biological neurons.

Network neuroscience adds another layer, using graph theory to analyze brain connectivity patterns and inspire new learning algorithms that respect how real neural networks are structured.

\hrulefill

\section*{Slide 11: Summary Table [10:45-11:45]}

\textcolor{DukeBlue}{\textbf{[GESTURE to table]}}

This table gives you a quick comparison of the key methods we've discussed.

\textcolor{DukeBlue}{\textbf{[TRACE rows as you speak]}} For neural decoding, EEGNet excels at generalizing across different BCI tasks, while motion-energy models enable that impressive visual reconstruction from fMRI we mentioned. For disease diagnosis, AlzheimerNet stands out with its ninety-eight point six seven percent accuracy, while YOLOv11 shows the power of combining multiple data types. For neuroscience-inspired AI, DNNs prove effective at modeling actual brain function, while Loihi demonstrates the energy efficiency advantage we can get from thinking more like biology.

\textcolor{red}{\textbf{[PAUSE, look at camera]}} The key takeaway? \ul{There's no one-size-fits-all solution}. Your choice of approach depends critically on your specific application, what kind of data you have, and your computational constraints. Sometimes a simple, efficient method beats a complex one.

\hrulefill

\section*{Slide 12: Key Results [11:45-12:15]}

\textcolor{DukeBlue}{\textbf{[GESTURE to slide]}}

This slide summarizes the impressive numbers we've already covered. \textcolor{DukeBlue}{\textbf{[POINT as you list]}} AlzheimerNet's \ul{ninety-eight point six seven percent accuracy} in detecting Alzheimer's stages. YOLOv11's successful multimodal fusion approach. EEGNet's superior BCI generalization across paradigms. And Loihi's stunning \ul{one thousand times better energy efficiency}.

\textcolor{red}{\textbf{[PAUSE, shift tone]}} These are real breakthroughs. But achievements alone don't tell the full story. Now let's talk about the critical challenges these systems still face.

\hrulefill

\section*{Slide 12: Challenges \& Ethics [12:15-13:45]}

\textcolor{DukeBlue}{\textbf{[SERIOUS, THOUGHTFUL TONE]}}

Despite all these achievements, we face significant challenges.

\textcolor{DukeBlue}{\textbf{[POINT left column, count]}} Four technical challenges stand out. \ul{First}, data complexity—going from single neurons to whole brain networks is like trying to understand a city by tracking every individual person. \ul{Second}, multimodal integration—each data type has different characteristics, noise profiles, and temporal resolutions. \ul{Third}, interpretability—the "black box" problem. \textcolor{green!60!black}{\textbf{[Example]}} A doctor won't use an AI that says "this patient has Alzheimer's" without explaining why. \ul{Fourth}, generalizability—brains are as unique as fingerprints, so what works for one person may fail for another.

\textcolor{red}{\textbf{[PAUSE]}}

\textcolor{DukeBlue}{\textbf{[POINT right column, shift to ethical tone]}} And then we have profound ethical questions. Privacy and autonomy: if we can decode neural signals, are your thoughts still private? \textcolor{green!60!black}{\textbf{[Example]}} Imagine an employer screening job candidates with brain scans. Cognitive enhancement: if AI can boost intelligence, who gets access? Will this create new forms of inequality? In clinical use: how do we ensure informed consent when patients don't understand how the AI works?

\textcolor{red}{\textbf{[PAUSE, look at camera]}} These aren't just philosophical questions—they're practical concerns we must address before widespread deployment.

\hrulefill

\section*{Slide 14: Future \& Priorities [13:45-15:15]}

\textcolor{DukeBlue}{\textbf{[SHIFT TO OPTIMISTIC TONE]}}

But I don't want to end on a somber note, because the future holds tremendous opportunities.

\textcolor{DukeBlue}{\textbf{[GESTURE left]}} Technically, we're moving toward real-time BCIs that can decode intentions instantly, closed-loop systems that adaptively respond to brain states—\textcolor{green!60!black}{\textbf{[Example]}} imagine a prosthetic arm that learns your movement preferences—and analysis of massive neural datasets that could reveal patterns invisible to human researchers. Neuromorphic computing promises AI that's both powerful and energy-efficient. Combining network neuroscience with machine learning could unlock how brain networks support cognition.

\textcolor{DukeBlue}{\textbf{[GESTURE right]}} Clinically, we're looking at neurorehabilitation that adapts to each patient, communication systems for people with locked-in syndrome, and entirely new treatments for neurological disorders.

\textcolor{red}{\textbf{[PAUSE]}}

But to get there, we need five research priorities. \ul{First}, explainable AI so clinicians trust the systems. \ul{Second}, standardized metrics so we can fairly compare approaches. \ul{Third}, better multimodal integration that maintains interpretability. \ul{Fourth}, personalized AI that adapts to individual differences. \ul{Fifth}, ethical frameworks developed through ongoing dialogue.

\textcolor{DukeBlue}{\textbf{[LOOK AT CAMERA]}} And all of this requires \ul{interdisciplinary collaboration}—computer scientists, neuroscientists, clinicians, and ethicists working together.

\hrulefill

\section*{Slide 15: Conclusion [15:15-16:30]}

\textcolor{DukeBlue}{\textbf{[CONFIDENT, WRAP-UP TONE]}}

Let me bring this together.

\textcolor{DukeBlue}{\textbf{[GESTURE left]}} My key contributions in this report: I've developed a taxonomy organizing AI applications in neuroscience, provided systematic analysis comparing methods across domains, identified the critical challenges holding us back, and examined the ethical questions we can't ignore.

The key takeaways? AI has \ul{fundamentally transformed} neuroscience—from how we collect data to how we understand the brain. The bidirectional relationship benefits both fields. We've seen remarkable achievements like ninety-eight percent accuracy disease detection and thousand-fold improvements in energy efficiency. But significant challenges remain in interpretability, generalizability, and ethics.

\textcolor{red}{\textbf{[PAUSE]}}

\textcolor{DukeBlue}{\textbf{[GESTURE right, optimistic]}} Looking forward, through interdisciplinary collaboration and careful ethical consideration, the integration of AI and neuroscience will benefit both fields and society. We're not just building better AI or understanding brains better—we're doing both simultaneously, and that's powerful.

\textcolor{red}{\textbf{[PAUSE]}}

\textcolor{DukeBlue}{\textbf{[SMILE, look at camera]}} This work was completed for COMPSCI 402 at Duke Kunshan University. I'm grateful to Professor Peng Sun for his guidance throughout this course.

\vspace{1em}

{\Large Thank you for your attention. I'm happy to answer questions.}

\textit{[Expected total time: 16:30 / Target: 16-18 minutes]}

\vspace{2em}

\hrulefill

\section*{Post-Recording Checklist}

\begin{itemize}
  \item[$\square$] Duration: 15-20 minutes
  \item[$\square$] Audio clear and audible
  \item[$\square$] Eye contact during key points (especially examples)
  \item[$\square$] Gestures used effectively
  \item[$\square$] Pace varies (slower for key numbers, faster for context)
  \item[$\square$] Enthusiasm evident throughout
  \item[$\square$] Examples made concepts tangible
  \item[$\square$] Transitions were smooth
  \item[$\square$] Slides clearly visible
\end{itemize}

\end{document}
